\section{Collisions in a regular conditional dynamics}\label{sec:collisions}

A real collision is a multiple zero of the signed readout, not a meeting
of eigenvalues of the fixed kinetic Hamiltonian. The state itself is
independent of which real observation parameter $k$ is queried.
Since $C(t,0)=1$, an event has $k_*\ne0$; evenness lets us discuss
$k_*>0$. The results of this section are conditional local descriptions,
not assertions that the canonical family has a collision.

\subsection{The probability and moment signature of a double event}
\begin{proposition}[Regular state, degenerate observation]\label{prop:event}
Suppose $0<t_*<T$ and
\begin{equation}\label{eq:double}
 C(t_*,k_*)=C_k(t_*,k_*)=0,\qquad A_*:=C_{kk}(t_*,k_*)\ne0.
\end{equation}
Then the success mass, state and count probabilities remain smooth, with
\begin{align}
 q_{0,+}=q_{0,-}&=p_*/2>0,&
 \partial_kq_{0,+}=\partial_kq_{0,-}&=0,\label{eq:eventprob}\\
 \partial_t(q_{0,+}-q_{0,-})&=\gamma p_*A_*\ne0,&
 d_1(t_*,k_*)&=-\gamma t_*p_*A_*\ne0.
 \label{eq:eventrate}
\end{align}
The event is equivalently described by the signed moments
\begin{equation}\label{eq:momenttriple}
 \E_*\cos(k_*X)=0,\qquad
 \E_*X\sin(k_*X)=0,\qquad
 \E_*X^2\cos(k_*X)=-A_*\ne0.
\end{equation}
\end{proposition}
\begin{proof}
Differentiate \eqref{eq:cosine} in $k$ using the finite moments.
This gives \eqref{eq:momenttriple}. At the event,
$(d_0)_{kk}=p_*A_*$ because $p_t$ is independent of $k$.
Equations~\eqref{eq:q}, \eqref{eq:heat} and \eqref{eq:dj} prove
\eqref{eq:eventprob}--\eqref{eq:eventrate}. Smoothness and the positive
floor for $p_t$ follow from Proposition~\ref{prop:state} and
\eqref{eq:selection}--\eqref{eq:speed}. The fixed-$j$ integrals give
smooth count probabilities. No division by the vanishing contrast is
needed for any state or probability formula.
\end{proof}

Both joint no-count outcome probabilities decrease by \eqref{eq:loss}.
At the event their values are equal, while their rates of decrease differ.
The one-count signed trace is exactly that difference multiplied by $-t_*$.
This is the dynamical content of the double-event curvature. It does not
make the event forbidden: the signed second moment has no sign constraint
from state positivity alone. Nor is $\mu_2$ a kinetic energy; its monotonicity
in \eqref{eq:tangent} is a property of the quadrature's second moment.

\subsection{Local zero geometry and the singularity of zero tracking}
\begin{theorem}[Local double-collision law]\label{thm:fold}
Under \eqref{eq:double}, let $a_*=(d_0)_{kk}(t_*,k_*)=p_*A_*$. Then
\begin{equation}\label{eq:fold}
 d_0(t_*+h,k_*+x)
 =\frac{a_*}{2}(x^2+2\gamma h)+O(|x|^3+|hx|+h^2).
\end{equation}
For $h<0$ sufficiently small there are two nearby real zeros,
\begin{equation}\label{eq:realroots}
 k_\pm(t)=k_*\pm\sqrt{-2\gamma h}+O(|h|).
\end{equation}
For $h>0$ the corresponding nearby analytic zeros form a conjugate pair,
\begin{equation}\label{eq:complexroots}
 z_\pm(t)=k_*\pm\ii\sqrt{2\gamma h}+O(|h|).
\end{equation}
All constants are local to the event. The complex roots refer only to
entire continuation of the source observation.
\end{theorem}
\begin{proof}
At the event $d_0=(d_0)_k=0$ and $(d_0)_t=\gamma a_*$ by
\eqref{eq:heat}. Finite multivariate Taylor expansion gives \eqref{eq:fold};
the required derivatives are paid by \eqref{eq:Hmixed}.
For $h=-s^2$ and $x=sy$, divide by $s^2$ and extend at $s=0$.
The limiting equation is $a_*(y^2/2-\gamma)=0$, with two simple real
roots. The analytic implicit-function theorem supplies real analytic
branches $y_\pm(s)=\pm\sqrt{2\gamma}+O(s)$, proving
\eqref{eq:realroots}. For $h=s^2$, the limiting roots are
$\pm\ii\sqrt{2\gamma}$, giving \eqref{eq:complexroots}.
A small fixed complex circle enclosing only the double zero at $h=0$
retains two zeros counted with multiplicity for small $h$, by uniform
analytic perturbation and the argument principle. The constructed roots
are therefore the full local pair. Reality of the source makes the
nonreal pair conjugate.
\end{proof}

Larger Newman parameter lies on the real-pair side because
$\lambda(t)-\lambda(t_*)=-\gamma h$. Reversing physical time reverses
the local picture. On a simple real zero curve $r(t)$,
\begin{equation}\label{eq:velocity}
 \dot r(t)=-\gamma\frac{(d_0)_{kk}}{(d_0)_k}
          =\frac{d_1}{t(d_0)_k}\quad(t>0).
\end{equation}
This follows by differentiating $d_0(t,r(t))=0$. Its divergence near the
double event belongs to locating the zero, while the state derivatives
remain finite. The squared gap has leading behavior
$(k_+-k_-)^2=-8\gamma h+o(|h|)$. Changing $\gamma$ rescales the clock;
it does not change an event's Newman parameter. The angular energies and
$\omega$ cancel from the co-moving observation, so they supply no additional
repulsion term in this local law.

For the original ternary observation $Y=+1,-1,0$ (the two successful
outcomes and failure),
\begin{equation}\label{eq:variance}
 \E Y=d_0,\qquad\Var Y=p_t-d_0^2.
\end{equation}
At the event this is $p_*>0$. A vanishing signed mean therefore has no
automatically vanishing single-trial noise. This exact variance calculation
is not a numerical experiment or a claim of optimal sample complexity.

\subsection{Balance of the positive and negative cosine sectors}
For fixed real $k\ne0$, let
$c_\pm(v)=\max(\pm\cos(kv),0)$ and define
\begin{equation}\label{eq:sectors}
 M_\pm=\E_t c_\pm(X)>0,\qquad
 B_\pm=\frac{\E_t[X^2c_\pm(X)]}{M_\pm},\qquad C=M_+-M_-.
\end{equation}
These are source-sector decompositions, distinct from the physical effects
$(I\pm\cos(kQ_t))/2$. Their boundaries are fixed when differentiating in
$t$, and finite moments justify the calculation despite the positive-part
corners. Equation~\eqref{eq:covariance} gives
\begin{equation}\label{eq:sectorratio}
 \partial_t\log(M_+/M_-)=-\gamma(B_+-B_-).
\end{equation}
At $C=0$, write $M_+=M_-=M$. Directly from the source integral,
\begin{equation}\label{eq:sectorcurvature}
 C_{kk}=-M(B_+-B_-),\qquad C_t=-\gamma M(B_+-B_-).
\end{equation}
Thus equal cosine masses, the stationary sine balance in
\eqref{eq:momenttriple}, and unequal sector second moments describe the
same double event. Quadratic attenuation changes the balance by suppressing
different radial locations at different rates. This interpretation and the
count-sector identity use the same signed source moments.

Cauchy--Schwarz gives the generic speed estimate
\begin{equation}\label{eq:covbound}
 |C_t|\le\gamma\sqrt{(\mu_4-\mu_2^2)
                   (\E_t\cos^2(kX)-C^2)}.
\end{equation}
It controls neither the sign nor the simultaneous vanishing of $C,C_k$.
It is available by direct source analysis and supplies no additional
collision margin merely through its dynamical interpretation.

\subsection{Higher multiplicities and temporal endpoints}
\begin{proposition}[Finite-order event polynomial]\label{prop:multiple}
If $d_0(t_*,\cdot)$ has exact finite order $m\ge2$ at $k_*$ and
$b_m=\partial_k^m d_0(t_*,k_*)\ne0$, then, uniformly for $(h,x)$ in a
fixed compact set,
\begin{equation}\label{eq:multiple}
 d_0(t_*+\delta^2h,k_*+\delta x)
 =\frac{b_m\delta^m}{m!}
 \sum_{r=0}^{\lfloor m/2\rfloor}
 \frac{m!\gamma^rh^rx^{m-2r}}{r!(m-2r)!}+o(\delta^m).
\end{equation}
For an interior event with $t_*>0$, if $m=2r$ then $d_j=0$ for $j<r$
and $d_r\ne0$. If $m=2r+1\ge3$, then $d_j=0$ for $j\le r$ and
$\partial_kd_r\ne0$.
\end{proposition}
\begin{proof}
Repeated differentiation of \eqref{eq:heat} gives
$\partial_t^a\partial_k^b d_0=\gamma^a\partial_k^{2a+b}d_0$.
In a finite multivariate Taylor expansion, every term of parabolic
weight $2a+b<m$ vanishes. Terms of weight $m$ give exactly
\eqref{eq:multiple}; a sufficiently high finite total-order remainder is
$o(\delta^m)$. The count assertions follow from \eqref{eq:dj} and the
nonzero factor $(-\gamma t_*)^j/j!$.
\end{proof}

The constants depend on $m$ and on the local event. A first-count diagnostic
therefore does not cover arbitrary multiplicities, and genericity is no
proof that the canonical family has only double events. At $t=0$ or $t=T$
only the physical side inside the slab is licensed. The entire source can
describe a two-sided germ without extending the apparatus beyond that slab.
The local laws are the familiar finite heat-flow geometry expressed in this
realization \cite{LocalCollisionGeometry,TNDDirectCollision}; no new
universal zero-collision law is claimed.
